\documentclass[11pt]{article}

\usepackage{amsmath, amssymb, amsthm}
\usepackage{geometry}
\usepackage{booktabs}
\usepackage{array}
\usepackage{xcolor}
\usepackage{titlesec}
\usepackage{parskip}

\geometry{margin=2.2cm}

\definecolor{accent}{RGB}{30,90,160}
\titleformat{\section}{\large\bfseries\color{accent}}{}{0em}{}[\titlerule]
\titleformat{\subsection}{\normalsize\bfseries}{}{0em}{}

\title{\textbf{Compute Core --- Mathematical Description}\\[0.4em]
       \large 5-Class Cardiac Arrhythmia Classifier, RBF-SVM ASIC}
\author{Adam Handwerger \quad \texttt{handwerg@pdx.edu}\\
        ECE410, Portland State University}
\date{2026-06-05}

\begin{document}
\maketitle
\noindent\rule{\linewidth}{0.5pt}

\vspace{0.3em}
The \texttt{svm\_compute\_core} implements two sequential subfunctions per support
vector: a \textbf{distance function} (\texttt{dist}) that computes the squared
Euclidean distance between an input feature vector and a support vector, and a
\textbf{kernel function} (\texttt{horner}) that evaluates the RBF exponential.
After all support vectors are processed, an \textbf{accumulation and classification}
step produces the output class label.

All arithmetic is performed in \(\mathbb{Q}_{6.10}\) 16-bit signed fixed-point
unless otherwise noted.

\vspace{0.5em}
\noindent\rule{\linewidth}{0.3pt}

% ─────────────────────────────────────────────────────────────────────────────
\section{Fixed-Point Representation (\(\mathbb{Q}_{6.10}\))}

A 16-bit signed value \(v\) is stored as the integer
\[
    \hat{v} = \left\lfloor v \cdot 2^{10} \right\rfloor \in [-32768,\; 32767]
\]
with the correspondence
\[
    v = \frac{\hat{v}}{1024}, \qquad
    \text{range: } [-32,\; +31.999], \qquad
    \text{LSB} = 2^{-10} \approx 0.001.
\]
Multiplication of two \(\mathbb{Q}_{6.10}\) values \(\hat{a}\) and \(\hat{b}\)
produces a 32-bit intermediate result that is right-shifted by 10 bits to
restore the \(\mathbb{Q}_{6.10}\) format:
\[
    \widehat{a \cdot b} = \frac{\hat{a} \cdot \hat{b}}{2^{10}}.
\]
The parameter \(\gamma = 0.25\) is exactly representable:
\(\hat{\gamma} = 0.25 \times 1024 = 256 = \texttt{0x0100}\).

% ─────────────────────────────────────────────────────────────────────────────
\section{Distance Function (\texttt{dist})}

\subsection{Definition}

For input feature vector \(\mathbf{x} \in \mathbb{Q}_{6.10}^{N}\) and support
vector \(\mathbf{sv} \in \mathbb{Q}_{6.10}^{N}\) with \(N = 256\) dimensions:
\[
    d^2(\mathbf{x},\, \mathbf{sv})
    \;=\;
    \sum_{i=0}^{N-1} \left( x_i - sv_i \right)^2.
\]

\subsection{Pipeline Implementation}

The accumulator is a two-stage registered pipeline:
\[
    \delta_i = x_i - sv_i,
    \qquad
    s_i = \delta_i^2,
    \qquad
    D_k = \sum_{i=0}^{k} s_i.
\]
Stage 1 (cycle \(c\)): compute and register \(s_i = \delta_i^2\).\\
Stage 2 (cycle \(c+1\)): add \(s_i\) to the running accumulator \(D\).

After the \(N\)-th input pair is presented, \textbf{two drain cycles} are
required before the result \(D_{N-1}\) is valid:
\begin{itemize}
    \item Drain cycle 1: \(s_{N-1}\) propagates from the multiplier register to
          the adder input.
    \item Drain cycle 2: \(s_{N-1}\) is added into \(D\) and the result is
          registered.
\end{itemize}
Total cycles per SV distance: \(N + 2 = 258\) at \(\text{LAT}=1\);
\quad \(N \cdot \text{LAT} + 2\) at general latency.

\subsection{Saturation}

The accumulator is 32-bit, with a hardware maximum of \(2^{32} - 1\).
The theoretical worst-case accumulated value for \(N=256\) and maximum
feature magnitude \(|x_i - sv_i| \leq 63.999\) is
\[
    D_{\max} = 256 \times (63.999)^2 \approx 1{,}048{,}544 \approx 2^{20},
\]
which is well within the 32-bit range. The fixed-point overflow boundary
is determined by the gamma multiply: after right-shifting by \(2^{10}\),
the maximum representable Q6.10 value from a 32-bit accumulator is
\[
    \frac{2^{32}}{2^{10}} = 2^{22}.
\]
With \(\hat{\gamma} = 2^8\), the accumulator would produce
\(\hat{p} > 2^{18}\) only if \(D > 2^{20}\), which does not occur on
real ECG data. The \texttt{ERR\_DIST\_OVERFLOW} flag is therefore a
defensive sanity check rather than a condition expected in normal
operation.

% ─────────────────────────────────────────────────────────────────────────────
\section{Kernel Function (\texttt{horner})}

\subsection{Definition}

The RBF (Gaussian) kernel is
\[
    K(\mathbf{x},\, \mathbf{sv})
    \;=\;
    e^{-\gamma \, d^2(\mathbf{x},\, \mathbf{sv})},
    \qquad \gamma = 0.25.
\]
The argument \(p = \gamma d^2\) ranges over \([0,\; \gamma D_{\max}] \approx [0,\; 262{,}136 / 1024] \approx [0,\; 256]\)
in floating-point. In \(\mathbb{Q}_{6.10}\) fixed-point, \(\hat{\gamma} = 256 = 2^8\)
and \(\hat{d}^2\) reaches up to \(\sim 2^{20}\) (saturating accumulator), so after
the \(\mathbb{Q}_{6.10}\) multiply (right-shift by \(2^{10}\)):
\[
    \hat{p} = \frac{\hat{\gamma} \cdot \hat{d}^2}{2^{10}}
    \leq \frac{2^8 \cdot 2^{20}}{2^{10}} = 2^{18}.
\]

\subsection{Range Reduction}

Direct Horner evaluation of \(e^{-p}\) over \([0, 256]\) requires degree
\(\geq 8\) and overflows \(\mathbb{Q}_{6.10}\) intermediate products.
Instead, \(p\) is split into integer and fractional parts:
\[
    I = \left\lfloor \frac{\hat{p}}{2^{10}} \right\rfloor \in [0,\; 255],
    \qquad
    F = \frac{\hat{p} \bmod 2^{10}}{2^{10}} \in [0,\; 1),
\]
so that
\[
    e^{-p} = e^{-I} \cdot e^{-F}.
\]

\subsubsection*{LUT term: \(e^{-I}\)}

A 256-entry read-only table stores
\[
    \texttt{lut}[I] = \left\lfloor e^{-I} \cdot 2^{10} \right\rfloor
    \in \mathbb{Q}_{6.10}, \qquad I = 0, 1, \ldots, 255.
\]
For \(I \geq 8\), \(e^{-I} < 0.00034\), which rounds to zero in
\(\mathbb{Q}_{6.10}\), so the effective range is \(I \in [0,\; 7]\).
The LUT occupies \(256 \times 16\) bits \(= 4\;\text{KB}\) of on-chip
flip-flop registers.

\subsubsection*{Polynomial term: \(e^{-F}\)}

For \(F \in [0,\; 1)\), numerical analysis shows that degree 6 is the minimum
Taylor polynomial that keeps the approximation error below the
\(\mathbb{Q}_{6.10}\) LSB of \(2^{-10} \approx 0.001\) across the full range.
The RTL defines up to 16 Taylor coefficients (\texttt{COEFF\_00}--\texttt{COEFF\_15})
and evaluates them via Horner's method. Defining signed coefficients
\(a_n = (-1)^n / n!\), the polynomial in standard form is:
\[
    e^{-F} \;\approx\; a_0 + a_1 F + a_2 F^2 + \cdots + a_d F^d,
\]
where:
\[
    a_0 = 1,\quad
    a_1 = -1,\quad
    a_2 = 0.5,\quad
    a_3 = -0.1\overline{6},\quad
    a_4 = 0.041\overline{6},\quad
    a_5 = -0.008\overline{3},\quad
    a_6 = 0.001\overline{388}.
\]
Rather than evaluating each power of \(F\) separately,
Horner's method computes this as a recurrence starting from the innermost term:
\[
    h_d = a_d, \qquad h_k = a_k + F \cdot h_{k+1}, \qquad k = d-1,\; d-2,\; \ldots,\; 0.
\]
The result is \(h_0 = e^{-F}\). The alternating signs are carried in the
\(a_n\) coefficients, so every step of the recurrence is a simple multiply-add
with no separate sign logic. Each step is one multiply and one add,
requiring \(d\) multiplications and \(d\) additions total.

\subsection{Kernel Output}

The two terms are combined by a \(\mathbb{Q}_{6.10}\) multiply:
\[
    \hat{K} = \frac{\texttt{lut}[I] \cdot \hat{e}^{-F}}{2^{10}}
    \;\in\; [0,\; 1024],
\]
where \(\hat{K} = 1024\) corresponds to \(K = 1.0\) (achieved when
\(d^2 = 0\), i.e.\ \(\mathbf{x} = \mathbf{sv}\)).

% ─────────────────────────────────────────────────────────────────────────────
\section{Alpha Accumulation and OVR Decision}

\subsection{Per-Class Accumulation}

For each support vector \(j = 0, \ldots, M-1\) with class label
\(\ell_j \in \{0,1,2,3,4\}\) and alpha coefficient
\(\alpha_j \in \mathbb{Q}_{6.10}\), the kernel output is weighted and
accumulated into the corresponding class score:
\[
    f_c(\mathbf{x}) \;=\;
    \sum_{\substack{j=0 \\ \ell_j = c}}^{M-1}
    \alpha_j \cdot K(\mathbf{x},\, \mathbf{sv}_j),
    \qquad c = 0, 1, 2, 3, 4.
\]
The five accumulators \(f_0, \ldots, f_4\) are 32-bit registers, updated
after each SV kernel evaluation. Alpha coefficients are stored on-chip in
a 500-entry register file (\texttt{alpha\_table}[]), loaded via the
Wishbone \texttt{ALPHA\_WR} register before classification begins.

\subsection{Classification (Argmax)}

After all \(M = 500\) support vectors are processed, the predicted class is
\[
    \hat{y} = \arg\max_{c \,\in\, \{0,1,2,3,4\}} f_c(\mathbf{x}).
\]
The argmax is implemented as a 5-way comparator tree and completes in one
clock cycle. The result is registered on \texttt{class\_out[2:0]} and
exposed on GPIO.

% ─────────────────────────────────────────────────────────────────────────────
\section{Operation Count Summary}

\begin{center}
\begin{tabular}{lrrr}
\toprule
\textbf{Stage} & \textbf{Operations per SV} & \textbf{SVs} & \textbf{Total} \\
\midrule
\texttt{dist}: subtractions & 256 & 500 & 128{,}000 \\
\texttt{dist}: multiplications & 256 & 500 & 128{,}000 \\
\texttt{dist}: accumulations & 256 & 500 & 128{,}000 \\
\texttt{horner}: multiplications & 3 & 500 & 1{,}500 \\
\texttt{horner}: additions & 3 & 500 & 1{,}500 \\
Alpha multiply-accumulate & 1 & 500 & 500 \\
\midrule
\textbf{Total per beat} & & & \textbf{387{,}500} \\
\bottomrule
\end{tabular}
\end{center}

\vspace{0.3em}
At \(N = 256\) and \(M = 500\), the design is strongly
\textbf{memory-bound}: arithmetic intensity \(AI \approx 1.5\)
OPs/byte (LUT preloaded), placing the operating point well below
the compute ceiling of the 40\;MHz fixed-point datapath.

% ─────────────────────────────────────────────────────────────────────────────
\section{FSM Cycle Budget (LAT = 3)}

\[
    T_{\text{beat}}
    = T_{\text{load}} + M \cdot \left( T_{\text{dist}} + T_{\text{kernel}} \right)
    + T_{\text{classify}}
\]
\[
    = N \cdot \text{LAT}
    + M \left( N \cdot \text{LAT} + 2 + 18 \right)
    + 1
\]
\[
    = 256 \times 3
    + 500 \times \left( 256 \times 3 + 20 \right)
    + 1
    = 768 + 500 \times 788 + 1
    = 394{,}769 \text{ cycles/beat}.
\]
At 40\;MHz: \(394{,}769 / 40 \times 10^{6} \approx 9.87\;\text{ms/beat}\),
consistent with the measured \(9.7\;\text{ms}\) (the small difference
reflects pipeline overlap between \textsc{load} and \textsc{compute}).

\vspace{1em}
\noindent\rule{\linewidth}{0.5pt}
{\small\itshape ECE410 $\cdot$ Portland State University $\cdot$
Adam Handwerger $\cdot$ sky130A $\cdot$ MIT-BIH Arrhythmia Database}

\end{document}
